\corrigesChapitre{5}{Géométrie et racines de l’unité}
\begin{corrige}{1}{Un quotient géométrique}
\begin{questions}
\item $\mathrm i$.
\item Module $1$, argument $\pi/2$ modulo $2\pi$.
\item $AB=AC$ et l’angle en $A$ est droit : triangle rectangle isocèle en $A$.
\end{questions}
\end{corrige}
\begin{corrige}{2}{Alignement}
\begin{questions}
\item $b-a=2+\mathrm i$, $c-a=4+2\mathrm i$.
\item $(c-a)/(b-a)=2$, réel positif.
\item Les points sont alignés et $B$ est le milieu de $[AC]$.
\end{questions}
\end{corrige}
\begin{corrige}{3}{Racines quatrièmes}
\begin{questions}
\item $1,\mathrm i,-1,-\mathrm i$.
\item Les quatre points sont aux intersections des axes et du cercle unité : ils forment un carré.
\item $1+\mathrm i-1-\mathrm i=0$.
\end{questions}
\end{corrige}
\begin{corrige}{4}{Racines cubiques}
\begin{questions}
\item $j=-1/2+\mathrm i\sqrt3/2$.
\item $j^2=-1/2-\mathrm i\sqrt3/2$ et $j^3=1$.
\item Les parties imaginaires se compensent et les parties réelles donnent $1-1/2-1/2=0$.
\end{questions}
\end{corrige}
\begin{corrige}{5}{Perpendicularité}
\begin{questions}
\item $(-2+4\mathrm i)/(2+\mathrm i)=2\mathrm i$.
\item Ce quotient est imaginaire pur non nul : les droites sont perpendiculaires.
\item Son module vaut $2$, donc $AC=2AB$.
\end{questions}
\end{corrige}
\begin{corrige}{6}{Un angle orienté}
\begin{questions}
\item $\arg b=\pi/4$, $\arg c=\pi/6$ modulo $2\pi$.
\item $\pi/6-\pi/4=-\pi/12$ modulo $2\pi$.
\item $AC/AB=2/\sqrt2=\sqrt2$.
\end{questions}
\end{corrige}
\begin{corrige}{7}{Un hexagone}
\begin{questions}
\item $\mathrm e^{\mathrm ik\pi/3}$ pour $k=0,\ldots,5$.
\item $2\sin(\pi/6)=1$ dans le cercle unité.
\item La somme vaut zéro, par somme géométrique ou par paires de sommets opposés.
\end{questions}
\end{corrige}
\begin{corrige}{8}{Un lieu d’angle}
\begin{questions}
\item $c/2=x/2+\mathrm iy/2$ est imaginaire pur si et seulement si $x=0$.
\item L’axe imaginaire privé de l’origine.
\item La distance $AC$ est libre ; seule la direction est imposée.
\end{questions}
\end{corrige}
\begin{corrige}{9}{Racines d’un autre complexe}
\begin{questions}
\item $|z|^3=8$, donc $|z|=2$.
\item $3\theta=\pi/2+2k\pi$, donc $\theta=\pi/6+2k\pi/3$.
\item $2\mathrm e^{\mathrm i\pi/6}$, $2\mathrm e^{5\mathrm i\pi/6}$, $2\mathrm e^{3\mathrm i\pi/2}$ ; il suffit de $k=0,1,2$.
\end{questions}
\end{corrige}
\begin{corrige}{10}{Spirale convergente}
\begin{questions}
\item $z_n=(\mathrm i/2)^n$.
\item $|z_n|=2^{-n}$ et un argument vaut $n\pi/2$.
\item Chaque étape tourne de $\pi/2$ et divise la distance à l’origine par $2$ ; les points tendent vers $O$.
\end{questions}
\end{corrige}
\begin{corrige}{11}{Avec un point fixe}
\begin{questions}
\item $\ell=1/(1-\mathrm i/2)=4/5+2\mathrm i/5$.
\item $z_{n+1}-\ell=(\mathrm i/2)(z_n-\ell)$, donc $z_n-\ell=-\ell(\mathrm i/2)^n$.
\item Le module de l’écart vaut $|\ell|2^{-n}\to0$ ; $z_n\to\ell$.
\end{questions}
\end{corrige}
\begin{corrige}{12}{La somme et le polygone}
\begin{questions}
\item $\omega\ne1$ et $\omega^n=1$ ; la somme géométrique vaut zéro.
\item L’affixe du point moyen est la somme divisée par $n$, donc zéro : le centre est le point moyen des sommets.
\item Non : le dénominateur $1-\omega$ s’annule et la somme vaut $1$. Il faut conserver l’hypothèse $n\geq2$ pour cette identité.
\end{questions}
\end{corrige}
